\paragraph{A quantitative logarithmic-square obstruction from the actual zero density.}
The preceding zero-gap argument proves that a coefficient
independent of Fourier height cannot work. A sharper
conclusion identifies the \emph{smallest possible leading
logarithmic-square growth rate} for a coefficient that
does depend on height. No RH assumption is used.

\begin{theorem}[Necessary scale of the spatial heat coefficient]
\label{thm:log-square-necessary}
Suppose $c_1(s,x)\ge0$ is finite at every real
$x$ and satisfies, at $s=0$,
\begin{equation}\label{eq:spatial-n1}
6U_2(0,x)\le c_1(0,x)U_1(0,x)
\qquad(x\in\mathbb R).
\end{equation}
Let $a>0$ range over the real zeros of $H_{1/2}$.
Then
\begin{equation}\label{eq:zeroexactenergy}
\boxed{\quad
c_1(0,a)\ge
6\sum_{\substack{b\in\mathbb R,\ H_{1/2}(b)=0\\b\ne a}}
\frac1{(a-b)^2}.\quad}
\end{equation}
Moreover, writing $N(R)$ for the number of zeros of
$H_{1/2}$ in $[R,2R]$, there exists such a zero
$a_R\in[R,2R]$ with
\begin{equation}\label{eq:finitezeroenergy}
c_1(0,a_R)\ge
\frac{12}{N(R)R^2}
\sum_{m=1}^{N(R)-1}
\frac{(N(R)-m)^3}{m^2}.
\end{equation}
Consequently
\begin{equation}\label{eq:logsquareobstruction}
\boxed{\quad
\limsup_{\substack{a\to+\infty\\H_{1/2}(a)=0}}
\frac{c_1(0,a)}{\log^2(2+a)}
\ge\frac18.\quad}
\end{equation}
In particular a bound
$c_1(0,x)\le C(1+\log^2(2+|x|))$
would require $C\ge1/8$; a bound
$c_1(0,x)=o(\log^2 x)$ is impossible.
\end{theorem}
\begin{proof}
At heat time $t=1/2$ all zeros of $H_t$ are
real. They are simple because this heat time lies
strictly above the de Bruijn--Newman constant,
for which the known upper bounds are less than
$1/2$ \cite{Polymath15}. Fix one zero $a$ and
factor $H_{1/2}(z)=(z-a)h(z)$ with $h(a)\ne0$.
The identities at $a$ are
\[
L_1[H_{1/2}](a)=h(a)^2,\qquad
L_2[H_{1/2}](a)
=h'(a)^2-h(a)h''(a).
\]
Since $H_{1/2}$ has an entire canonical product
of genus at most one, its logarithmic second
derivative at $a$ yields
\[
\frac{L_2[H_{1/2}](a)}
{L_1[H_{1/2}](a)}
=-(\log h)''(a)
=\sum_{\substack{b\ne a\\H_{1/2}(b)=0}}
\frac1{(a-b)^2}.
\]
Equation \eqref{eq:spatial-n1} now proves
\eqref{eq:zeroexactenergy}.

Let $a_1<\cdots<a_N$ be all zeros in $[R,2R]$
and $N=N(R)\ge2$. For each positive integer
$m<N$ set
$d_{i,m}=a_{i+m}-a_i>0$.
Every consecutive gap occurs in at most
$m$ such differences, so
$\sum_{i=1}^{N-m}d_{i,m}\le mR$.
Convexity of $y\mapsto y^{-2}$, or
Jensen's inequality, gives the finite
exact estimate
\[
\sum_{i=1}^{N-m}d_{i,m}^{-2}
\ge\frac{(N-m)^3}{
 (\sum_i d_{i,m})^2}
\ge\frac{(N-m)^3}{m^2R^2}.
\]
Each unordered pair contributes twice to
the sum of the energies in
\eqref{eq:zeroexactenergy}. Averaging over
the $N$ zeros therefore yields
\[
\max_{1\le i\le N}c_1(0,a_i)
\ge\frac{12}{NR^2}
\sum_{m=1}^{N-1}\frac{(N-m)^3}{m^2},
\]
which is \eqref{eq:finitezeroenergy}.

The established zero-counting asymptotic for
the real $H_{1/2}$ zeros is
\[
N(R)=\frac{R\log R}{4\pi}(1+o(1))
\qquad(R\to\infty);
\]
see the classical de Bruijn--Newman zero-count
analysis \cite{RodgersTao,Polymath15}.
For each fixed $m$, $(N-m)^3/N^3\to1$,
and the sum over $m$ is bounded by
$\sum_{m\ge1}m^{-2}=\pi^2/6$.
Dominated convergence therefore gives
\[
\frac1{N^3}\sum_{m=1}^{N-1}
\frac{(N-m)^3}{m^2}
\longrightarrow\frac{\pi^2}{6}.
\]
Substituting into \eqref{eq:finitezeroenergy},
using $\log(2+a_R)\sim\log R$, proves
\eqref{eq:logsquareobstruction}.
\end{proof}

The coefficient $1/8$ is the inverse-square
energy of the asymptotically equally spaced
zero lattice in this Fourier normalization;
it is a necessary threshold, \emph{not}
a verified upper bound for the genuine theta
Laguerre quotient. The argument strengthens the
zero-density obstruction without solving it:
a candidate logarithmic-square envelope must
still be proved simultaneously for every
spatial height, heat time and Laguerre order.
The existing near-collision obstruction is
also untouched.

The elementary inverse-gap inequalities are
tested by the exact-rational checker
\path{certificates/xi_heat_log_square_gap_exact.py}
and their two-gap polynomial gate is verified by
\path{XiKernel/LogSquareHeatGap.lean}.
The canonical-product identity, the zero-count
asymptotic and Jensen averaging are manuscript
arguments, not end-to-end Lean results.
